\section{ZAKLJUČAK}
Autonomni viličari i industrijska vozila predstavljaju ključni stup tranzicije prema konceptu pametnih skladišta i šire vizije *Industrije 4.0 \cite{kion}*. Tehnologija je prešla dug put od Barrettovog prvog žicom vođenog tegljača iz 1953. godine do današnjih visokosofisticiranih autonomnih viličara koji koriste SLAM algoritme, 3D LiDAR-e i računalni vid.

Glavni pokretač ove tranzicije je potreba za fleksibilnošću. Prirodna navigacija omogućila je uvođenje automatizacije u već postojeća skladišta bez potrebe za skupim građevinskim zahvatima i postavljenjem fizičkih vodiča. Sustavi fuzije senzora uspješno spajaju podatke s LiDAR-a, IMU-a, odometrije i kamera kako bi osigurali milimetarsku preciznost manipulacije teretom te maksimalnu sigurnost pješaka u skladu sa standardom ISO 3691-4.

Iako visoki inicijalni troškovi implementacije i održavanja te osjetljivost na vanjske uvjete predstavljaju izazove za mala i srednja poduzeća, prednosti u vidu 24-satnog rada, dramatičnog smanjenja oštećenja robe i optimizacije procesa čine ovu tehnologiju ekonomski isplativom na srednji i dugi rok. Budući razvoj intralogistike bit će obilježen daljnjom integracijom umjetne inteligencije za prediktivnu analizu tokova materijala, uvođenjem 5G mreža za bržu komunikaciju flote te razvojem hibridnih kolaborativnih sustava u kojima ljudi i autonomni strojevi rade u sinergiji.
